\section{Research targets, results, and source status}
\label{sec:scope}

The question motivating this continuation is how much of a singular
\emph{ordered} nested sum can be evaluated using the familiar Gamma,
polygamma, Hurwitz-zeta and Stieltjes functions. Full symmetrization
removes ordering information and often reduces by the stuffle product.
The individual ordered germ retains that information. Its singularities
also meet along several hyperplanes, so a constant Laurent coefficient
must specify the approach to their intersection.

We work from the repository snapshot
\begin{center}
\small\texttt{4c1286161c425eca957b4b9f7cd436c37a4c98f3}.
\end{center}
The source survey includes the canonical manuscript, its status records,
and the seven ZIP packages then present in \src{docs/incoming}. The
mathematical audit concentrates on the relevant ordered-germ,
coincident-product, resonance, and Gaussian-status material; it is
not a re-verification of every proof in the full manuscript.
The exact files and archive hashes are recorded in the accompanying
\src{integration/source_snapshot.json}. All sources were inspected as
research evidence; successful numerical replay is not treated as a
proof of an analytic theorem.

\subsection{The concrete questions answered}

The incoming \emph{Resonant Gamma Jets and Directional Zeta Identities}
report asks for an ordered depth-three germ with independent exponents,
its intersecting polar divisors, a normally convergent remainder, and
the constant along every transverse ray. It separately asks for
tangential curves and the exact number of curve coefficients needed
to determine a finite part \cite[Further research, Q1 and Q3]{repojets}.
Both questions are answered here. The first answer also yields the
following all-depth specialization:
\begin{equation}\label{scope:mainfamily}
 \FP_{\eps=0}Z_d(1+p\eps,1+q\eps,\ldots,1+q\eps;a),
 \qquad a>0.
\end{equation}
It closes in ordinary Hurwitz-zeta derivatives and Stieltjes constants
whenever all prefix slopes $p+jq$ are nonzero. The initial slope may
differ from every trailing slope; no symmetrization of the ordered sum
is needed.

The \emph{Coincident Stieltjes} report asks for a cubic closure theorem
and specifically for the coordinate-normalized digamma cube
\cite[Further research: cubic coincident moments]{repocoincident}.
We construct the full cubic completion and evaluate that finite part
in terms of a single precisely defined diagonal Tornheim derivative.
This is a finite reduction to an established higher-depth function.
A further reduction of that derivative to ordinary constants is a
separate unanswered question.

\begin{center}
\begin{tabular}{>{\raggedright\arraybackslash}p{.26\textwidth}>{\raggedright\arraybackslash}p{.64\textwidth}}
\toprule
Result & Exact outcome\\
\midrule
Ordered local germs & Convergent remainder construction at every depth and
every common positive shift; finite coefficient recursion.\\[4pt]
One distinguished slope & Every nonpositive Laurent coefficient reduces
to finitely many ordinary Hurwitz and Stieltjes jets.\\[4pt]
Independent depth three & Complete pole part and constant term; one
harmonic Stieltjes coefficient contains the residual ordering dependence.\\[4pt]
Tangential curves & Pole order $\sum\nu_j$ and a sharp required curve
jet order $\sum\nu_j+\max\nu_j$.\\[4pt]
Cubic coincident products & All Stieltjes indices and all argument
derivatives; explicit digamma-cube reduction to $\tau'''(0)$.\\
\bottomrule
\end{tabular}
\end{center}

\subsection{What is classical and what is added to this programme}

Laurent expansions of Euler--Zagier multiple zeta functions have a
substantial literature. Matsumoto, Onozuka and Wakabayashi analyze Laurent
coefficients at integer points \cite{MOW}. Saha constructs regularized
multiple Stieltjes germs and, in particular, gives the canonical
ordered decomposition at the all-one point \cite[Theorem 2]{Saha}.
Our tail-defect proof realizes that structure for an arbitrary
positive common Hurwitz shift. The explicit slope evaluations, the
reduction of every transverse depth-three directional finite part to
one named harmonic coefficient and ordinary Hurwitz data, and the
sharp curve-jet statement are the
specific conclusions developed for the repository's current questions.

The harmonic Hurwitz-zeta convention is that of Karg\i n, Dil,
Cenkci and Can \cite{Kargin}; the strict ordered convention used
here differs from their inclusive sum by a single Hurwitz-zeta term.
We give the conversion rather than silently identifying the constants.
The triple-Hurwitz Fourier relation is part of the Tornheim literature
of Espinosa and Moll \cite{EMTornheim}. The Mellin method follows the
established approach discussed by Borwein and Dilcher and by Bailey
and Borwein \cite{BorweinDilcher,BaileyBorwein}. Low diagonal derivatives
are recovered and used as consistency checks, without claiming their
first evaluation. The bibliography is a targeted source comparison,
not an exhaustive proof of publication priority.

\subsection{Conventions and mathematical status}

For $a>0$ we use real logarithms of $n+a$, and
\begin{equation}\label{scope:stieltjes}
 \zeta(1+u,a)=\frac1u+
  \sum_{m\ge0}\frac{(-1)^m}{m!}\gamma_m(a)u^m,
 \qquad \gamma_0(a)=-\psi(a).
\end{equation}
A prime on $\zeta$ denotes its first, spectral derivative. For $m\ge2$,
\begin{equation}\label{scope:polygamma}
 \zeta(m,a)=\frac{(-1)^m}{(m-1)!}\psi^{(m-1)}(a).
\end{equation}
These conventions follow the standard Hurwitz calculus \cite{DLMFHurwitz}.
Without a shift, $\gamma_m=\gamma_m(1)$ and $\gamma=\gamma_0(1)$.
The rising factorial is $(z)_k=\prod_{j=0}^{k-1}(z+j)$, with
$(z)_0=1$.
Brackets $[u^j]$ or $[y^jt^k]$ mean ordinary Taylor coefficients,
so conversion to derivatives requires the displayed factorials.

For a specified meromorphic function of one variable, $\FP$ means its
coefficient of $\eps^0$. For a spatial endpoint integral, it means
the constant remaining in the cutoff expansion in the stated spatial
coordinate. These two operations are related by explicit corrections
in Section~\ref{sec:cubic}; they are not interchangeable definitions.
The holomorphic regular germ is a third object, defined by an exact
subtraction recursion.

The canonical manuscript already proves $S_4$ and retains the short
$S_6$ and revised $S_8$ candidates as conjectures \cite{repomain}.
No argument in this article changes that status. We do not infer
period independence, minimal arithmetic bases, or nonreducibility
from a failed search. Every theorem below is proved analytically;
the attached computations check finite algebra and independently
evaluated examples.
